\chapter{Simulation contre production}\label{ch:rs:simulation-versus-live}

Le rejeu de recherche d'une stratégie de cotation, à savoir la simulation de niveau 3 du chapitre~18 lancée sur les mêmes heures que celles où la stratégie a ensuite opéré, annonçait une perte
d'environ 50~\$ de l'heure. En production, elle a perdu 491~\$ de l'heure, et a exécuté deux fois plus de lots. Le code du rejeu n'avait rien de faux~; ses exécutions avaient été reproduites
dans trois langages. Ce qui a mal tourné, c'est ce que rencontre toute société lors de son premier mois en production~: le simulateur répondait à une question légèrement différente de celle que posait
le marché. Ce chapitre construit le rapprochement (\texttt{firm.simlive})~: apparier les exécutions réelles aux exécutions simulées, faire passer le P\&L simulé au
P\&L réel, une hypothèse à la fois, repérer les exécutions que le simulateur ne pouvait pas produire, et décider de ce qu'il faut modifier pour que l'estimation suivante soit plus proche. Le
«~marché de production~» est un substitut que l'on peut rejouer~: \texttt{firm.tape}, avec la stratégie qui opère à l'intérieur comme agent, ses ordres étant de vrais ordres que d'autres traders
voient et contre lesquels ils traitent.

\section{Rapprochement des exécutions}

\begin{definition}[Rapprochement simulation--production, trading fantôme]\label{def:rs:simulation-versus-live:reconciliation}
Le \emph{rapprochement simulation--production}\index{rapprochement simulation--production} (sim-to-live reconciliation) est la comparaison systématique des ordres, des exécutions et du P\&L réels d'une stratégie avec ce que produit son
simulateur sur les mêmes données de marché et les mêmes signaux. Le \emph{trading fantôme}\index{trading fantôme} (shadow trading) fait tourner le simulateur en parallèle du trading réel,
sur le flux de cotations en direct, de sorte que la comparaison soit disponible chaque jour.
\end{definition}

Le quoteur du chapitre conserve un lot sur le meilleur cours acheteur et un lot sur le meilleur cours vendeur, dans la limite de cinq lots d'inventaire. En production, à l'intérieur de \texttt{firm.tape}, il tire ses propres latences
(données de marché 0.02 seconde plus une exponentielle de moyenne 0.02~; saisie d'ordre 0.03 plus une exponentielle de moyenne 0.02) et paie 0.1 tick par action. Ses ordres
sont appariés selon la priorité temporelle, comptent dans le carnet d'ordres auquel réagissent les traders informés et les fournisseurs de liquidité, et ne sont jamais annulés par le flux de fond. Le
rejeu de l'estimation de recherche est le rejeu du chapitre~18 (modèle de file FIFO) sur la même heure, sans le quoteur~: la même graine et le même chemin de prix efficient, comme
l'aurait lancé un chercheur la veille. Sur six sessions d'une heure, le quoteur a exécuté 597 lots par heure en production~; le rejeu, sur les messages de l'heure réelle, sans
ceux du quoteur et avec les latences réelles, en a exécuté 305.

\begin{definition}[Test de parité du rejeu]\label{def:rs:simulation-versus-live:parity}
Un \emph{test de parité du rejeu}\index{test de parité du rejeu} (replay parity test) rejoue une session réelle dans le simulateur, avec les décisions et les instants des ordres réels, et mesure la
part des exécutions réelles que le simulateur reproduit (même sens, même prix et même instant à une tolérance près) et la part des exécutions simulées qui ont eu lieu en production.
\end{definition}

\texttt{firm.simlive.match\_fills} associe chaque exécution réelle à une exécution simulée de même sens et de même prix à une seconde près. Seuls 27~\% des lots réels ont une
contrepartie simulée, et 54~\% des lots simulés ont eu lieu en production. Une parité aussi faible indique que le simulateur n'est pas une version bruitée de la production~: c'est un autre modèle de celle-ci.

\section{Rapprochement du P\&L}

\begin{definition}[Décomposition des écarts]\label{def:rs:simulation-versus-live:decomposition}
Une \emph{décomposition des écarts}\index{décomposition des écarts} (divergence decomposition) fait passer le P\&L simulé au P\&L réel par une suite d'estimations intermédiaires, chacune remplaçant
une hypothèse par sa valeur réelle, et attribue la variation de chaque étape à cette hypothèse.
\end{definition}

L'ordre des étapes compte (changer d'abord la latence puis le marché n'équivaut pas à l'inverse), et chaque étape doit être un passage du simulateur, non une
estimation. La suite du chapitre, moyennée sur les six sessions (la \cref{fig:rs:simulation-versus-live:waterfall})~:
\begin{enumerate}
\item le rejeu de recherche~: ni latence, ni frais, l'heure sans le quoteur~: $-\$50.5$ de l'heure~;
\item avec les latences moyennes réelles~: $-\$56.4$ (latence~: $-\$5.9$)~;
\item sur le marché que le quoteur a effectivement rencontré, les messages de l'heure réelle privés de ses propres messages~: $-\$112.8$ (le marché~: $-\$56.4$)~;
\item les exécutions réelles elles-mêmes, avant frais~: $-\$431.2$ (les exécutions~: $-\$318.3$)~;
\item en production, après frais~: $-\$490.9$ (frais~: $-\$59.7$).
\end{enumerate}

L'étape «~marché~» compare deux réalisations différentes du flux d'ordres de l'heure, avec et sans le quoteur, et une part en relève du hasard. Une seconde réalisation de chaque heure
sans le quoteur (même chemin de prix efficient, autre graine de flux d'ordres) diffère de la première de 62.8~\$ de l'heure en moyenne~: l'étape «~marché~» se situe dans la marge du hasard.
La latence et les frais sont faibles et connus~; l'écart se trouve dans les exécutions.

\begin{omfigure}
\begin{tikzpicture}
\begin{axis}[width=11cm,height=5.8cm,ybar,bar width=20pt,ylabel={P\&L de l'heure (\$)},tick label style={font=\small},label style={font=\small},
  xtick={0,1,2,3,4},xticklabels={{rejeu de\\recherche},{+ latence\\réelle},{+ marché\\rencontré},{+ exécutions\\réelles},{+ frais~:\\production}},
  xticklabel style={align=center,font=\footnotesize},xmin=-0.6,xmax=4.6,ymin=-560,ymax=0,grid=major,grid style={gray!25},
  nodes near coords,nodes near coords style={font=\footnotesize,/pgf/number format/fixed,/pgf/number format/fixed zerofill,/pgf/number format/precision=1},
  table/col sep=comma]
\addplot[fill=omThm!45,draw=omThm] table[x=k,y=level]{figdata/research/19-simulation-versus-live/waterfall.csv};
\end{axis}
\end{tikzpicture}

\omcaption{Du rejeu de recherche à la production, une hypothèse à la fois~: le P\&L de l'heure du quoteur, moyenné sur six sessions simulées, à chaque étape de la
décomposition, chaque barre donnant le niveau atteint après cette étape. Données~: \texttt{rs\_simlive.decomposition}.}\label{fig:rs:simulation-versus-live:waterfall}
\end{omfigure}

\section{Localiser les écarts}

Deux mécanismes produisent les exécutions réelles que le rejeu ne peut pas reproduire. Ils ne sont visibles que dans les messages réels, où figurent les ordres propres du quoteur.

\textbf{Transactions arrêtées au quoteur.} Un ordre agressif qui rencontre le lot du quoteur en tête de file est exécuté par celui-ci et peut s'arrêter là. Dans les
messages réels privés de ceux du quoteur, cette transaction n'existe pas~; un rejeu de ceux-ci ne peut pas remplir le quoteur simulé avec un volume qu'il ne voit jamais. 68~\% des lots réels proviennent d'ordres agressifs
qui ne se sont exécutés que contre le quoteur. Le marché après qu'une société a opéré n'est pas le marché sans cette société.

\textbf{Le dernier ordre à un prix.} Lorsque les autres ordres à ce prix sont annulés à mesure que le prix s'éloigne, le quoteur, dont l'annulation est encore en route vers la
bourse, reste seul à un prix périmé, et l'agresseur suivant le prend. Dans le rejeu, le niveau se vide simplement, et rien ne s'y échange. 21~\% des lots réels ont été
exécutés sans autre ordre au repos à leur prix, et leur mark-out après dix secondes était de $-0.85$ ticks, contre $-0.11$ pour les autres exécutions réelles~: les pires exécutions
de la journée sont celles que le simulateur ne montre jamais.

\section{Garder le simulateur fidèle}

\begin{definition}[Calibration du simulateur, implementation shortfall, prix d'arrivée]\label{def:rs:simulation-versus-live:calibration}
La \emph{calibration du simulateur}\index{calibration du simulateur} (simulator calibration) ajuste les paramètres d'un simulateur (latence, modèle de file, probabilités d'exécution) de sorte que ses résultats coïncident avec les
résultats observés en production. L'\emph{implementation shortfall}\index{implementation shortfall} d'un ordre (Perold, 1988) est l'écart entre le P\&L d'un portefeuille fictif négocié
en totalité au prix de l'\omterm{def:rs:vectorised-backtests:timing}{instant de décision}, et le P\&L effectivement réalisé~; le prix à l'\omterm{def:rs:vectorised-backtests:timing}{instant de décision}, ou à l'arrivée de l'ordre sur le marché, est son \emph{prix
d'arrivée}\index{prix d'arrivée}.
\end{definition}

La calibration est la réponse évidente, et ici elle échoue de façon instructive. Aucune latence ne permet au rejeu d'exécuter autant de lots que la production~: même à latence nulle, le rejeu sur les
messages réels exécute 1\,898 lots sur les six sessions, contre 3\,582 en production. Le passage au modèle de file en tête exécute 8\,886 lots, davantage qu'en production, avec un P\&L de $+\$2\,207$
contre $-\$2\,587$ en production. Un simulateur calé pour égaler le nombre d'exécutions aurait triplé l'erreur sur le P\&L. La divergence est structurelle, et le remède l'est aussi~: une
simulation réactive, dans laquelle les ordres de la stratégie font partie du marché (comme agent ici, et comme simulateur de bourse du livre~10), pour les stratégies dont les ordres
représentent une part notable du flux à leurs prix~; et pour les autres, un rejeu tenu fidèle par les chiffres que ce chapitre produit chaque jour.

Le rapprochement permanent est un rapport, produit chaque jour à partir du \omterm{def:rs:simulation-versus-live:reconciliation}{trading fantôme}~: parité des exécutions (les deux parts), étapes de la décomposition, les mark-out des exécutions réelles par
groupe (appariées, réelles seulement, seules à leur prix), et l'\omterm{def:rs:simulation-versus-live:calibration}{implementation shortfall} de chaque ordre parent par rapport à son \omterm{def:rs:simulation-versus-live:calibration}{prix d'arrivée}, réparti entre exécution, coût d'opportunité et
frais (\texttt{firm.simlive.implementation\_shortfall}). Un simulateur qui n'a pas été rapproché de la production est une hypothèse~; un simulateur rapproché chaque jour est un
instrument.

\begin{omfigure}
\begin{tikzpicture}
\begin{axis}[width=11cm,height=5.4cm,ybar,bar width=8pt,xlabel={session},ylabel={P\&L (\$ de l'heure)},tick label style={font=\small},label style={font=\small},
  xtick={1,2,3,4,5,6},ymin=-1050,ymax=150,grid=major,grid style={gray!25},legend columns=3,
  legend style={font=\footnotesize,at={(0.5,-0.3)},anchor=north,draw=none,fill=none,/tikz/every even column/.append style={column sep=8pt}},
  table/col sep=comma]
\addplot[fill=omDef!55,draw=omDef] table[x=session,y=sim]{figdata/research/19-simulation-versus-live/sessions.csv};
\addplot[fill=gray!40,draw=gray] table[x=session,y=chance]{figdata/research/19-simulation-versus-live/sessions.csv};
\addplot[fill=omThm!55,draw=omThm] table[x=session,y=live]{figdata/research/19-simulation-versus-live/sessions.csv};
\legend{rejeu de recherche,autre réalisation du rejeu,production}
\end{axis}
\end{tikzpicture}

\omcaption{Le P\&L du teneur de cotation dans chacune des six sessions simulées~: le rejeu de recherche, le rejeu d'une seconde réalisation de la même heure (avec les latences de production) et la production.
La perte en production dépasse les deux autres dans chaque session. Données~: \texttt{rs\_simlive.decomposition}.}\label{fig:rs:simulation-versus-live:sessions}
\end{omfigure}

\section{Tutoriel~: dix fois la perte simulée}\label{tut:rs:simulation-versus-live:tut}

\begin{tutorial}
\textbf{Objectif.} Faire trader le teneur de cotation en production dans \texttt{firm.tape}, rejouer les mêmes heures et procéder au rapprochement~: parité, décomposition, localisation de l'écart et une
calibration qui échoue. \textbf{État final~:} \cref{fig:rs:simulation-versus-live:waterfall,fig:rs:simulation-versus-live:sessions}~; les chiffres de ce chapitre.
\begin{tutsteps}
\item \textbf{Apparier les exécutions}~: glouton, dans l'ordre temporel, par sens et par prix à une tolérance près.
\omcode{code/firm/simlive/firm_simlive.py}{31}{51}{Appariement des exécutions de production avec les exécutions simulées.}\label{lst:rs:simulation-versus-live:match}
\item \textbf{L'agent}~: le teneur de cotation au sein du marché, avec ses propres latences.
\omcode{code/research/19-simulation-versus-live/python/rs_simlive.py}{36}{70}{Le teneur de cotation en production, en tant qu'agent de firm.tape.}\label{lst:rs:simulation-versus-live:agent}
\item \textbf{Exécuter} \texttt{decomposition()}, \texttt{steps()}, \texttt{vanished\_volume()}, \texttt{alone\_at\_price()}, \texttt{effective\_latency()},
\texttt{model\_calibration()} et \texttt{fig\_simlive.py}.
\end{tutsteps}
\textbf{Pour aller plus loin.} Faire annuler le teneur de cotation dès que les autres ordres à son prix passent sous un seuil, et observer les exécutions «~seules~» et la réduction de l'écart~;
rétablir, dans les données d'entrée du rejeu, les exécutions contre le teneur de cotation comme des exécutions contre un ordre anonyme, puis mesurer à nouveau la parité.
\end{tutorial}

\section{Construction~: la boîte à outils du rapprochement}\label{bld:rs:simulation-versus-live:simlive}

\begin{build}
\bfield{Rôle} Comparaison quotidienne de chaque stratégie en production avec son simulateur, et éléments de preuve pour modifier le simulateur.
\bfield{Interface} \texttt{FILL}, \texttt{as\_fills}, \texttt{match\_fills(live, sim, tol)}, \texttt{parity}, \texttt{fill\_pnl(fills, mark)}, \texttt{waterfall(steps)},
\texttt{calibrate(grid, metric, target)}, \texttt{implementation\_shortfall(side, qty, decision\_px, fills, end\_px, fee\_per\_unit)}~; \texttt{firm.tape.simulate(cfg,
agent)} pour le substitut réactif.
\bfield{Règles} Chaque étape d'une décomposition est un passage, dans un ordre énoncé~; la parité est indiquée dans les deux sens~; la calibration se juge sur le P\&L et les mark-outs, et non sur les seuls nombres d'exécutions.
\bfield{Tests d'acceptation} \texttt{code/firm/simlive/tests/}~: appariement avec quantités partielles et une tolérance~; parité dans les deux sens, et parfaite sur des exécutions identiques~; P\&L par exécution
sommant au total valorisé~; les étapes de la cascade~; calibration sur une grille~; l'\omterm{def:rs:simulation-versus-live:calibration}{implementation shortfall} d'un achat et d'une vente, calculée à la main.
\bfield{Pour aller plus loin} Appariement par identifiant d'ordre avec les accusés de réception de la bourse~; décomposition par heure de la journée~; recalibration d'un modèle probabiliste de file d'attente sur les exécutions de production.
\end{build}

\begin{omsources}
\item A.~F.~Perold, ``The implementation shortfall: paper versus reality'', \emph{Journal of Portfolio Management} 14(3), 1988.
\end{omsources}

\section{Exercices}

\begin{exercise}[$\star$]\label{exo:rs:simulation-versus-live:1}
Exécutions en production~: achat de 100 à 50.00 à 10.2~s, vente de 100 à 50.02 à 14.0~s. Exécutions simulées~: achat de 100 à 50.00 à 10.9~s, achat de 100 à 49.99 à 12.0~s. Avec une
tolérance d'une seconde, quelles sont les deux parts de parité~?
\end{exercise}

\begin{exercise}[$\star$]\label{exo:rs:simulation-versus-live:2}
Une décision d'achat de 10\,000 actions est prise à 20.00. 6\,000 sont exécutées à un prix moyen de 20.04 et le reste est annulé~; le cours termine à 20.30. Les frais sont de 0.2 cent par action.
Calculer l'\omterm{def:rs:simulation-versus-live:calibration}{implementation shortfall} et ses trois composantes.
\end{exercise}

\begin{exercise}[$\star$]\label{exo:rs:simulation-versus-live:3}
Pourquoi chaque étape d'une \omterm{def:rs:simulation-versus-live:decomposition}{décomposition des écarts} doit-elle être un passage du simulateur plutôt qu'une estimation~?
\end{exercise}

\begin{exercise}[$\star\star$]\label{exo:rs:simulation-versus-live:4}
Pourquoi l'étape de marché de la décomposition peut-elle rester dans les limites du hasard alors que l'impact de la stratégie est réel~? Comment séparer les deux~?
\end{exercise}

\begin{exercise}[$\star\star$]\label{exo:rs:simulation-versus-live:5}
Expliquer pourquoi un rejeu de la session de production sans les messages propres à la stratégie sous-estime ses exécutions, et quelle part des lots de production du chapitre il ne peut pas voir.
\end{exercise}

\begin{exercise}[$\star\star$]\label{exo:rs:simulation-versus-live:6}
Pourquoi les exécutions qui ont lieu lorsque la stratégie est le dernier ordre à son prix sont-elles moins bonnes que les autres~?
\end{exercise}

\begin{exercise}[$\star\star\star$]\label{exo:rs:simulation-versus-live:7}
\emph{Programmation.} Avec \texttt{rs\_simlive.effective\_latency}, vérifier qu'aucune latence ne comble l'écart d'exécutions. Lequel des deux mécanismes du chapitre une modification de latence
affecterait-elle, et pourquoi pas l'autre~?
\end{exercise}

\begin{exercise}[$\star\star\star$]\label{exo:rs:simulation-versus-live:8}
\emph{Trouver la faille.} «~Les exécutions en production étaient le double de celles du simulateur~; nous avons donc basculé le simulateur vers des exécutions optimistes~; désormais les nombres d'exécutions coïncident et le \omterm{def:rs:vectorised-backtests:backtest}{backtest} est
calibré.~»
\end{exercise}

\section{Problème~: dix fois la perte simulée}

\begin{problem}[{Problème du week-end --- une stratégie en production face à son simulateur}]\label{pb:rs:simulation-versus-live:1}
Le teneur de cotation du chapitre, en production dans \texttt{firm.tape} pendant six sessions d'une heure, et ses rejeux de niveau 3.

\medskip\noindent\textbf{Partie I --- Exécutions.}
\begin{enumerate}
\item Combien de lots par heure le teneur de cotation a-t-il exécutés en production, et combien le rejeu en a-t-il exécutés sur les messages de production sans les siens~?
\item Quelles sont les deux parts de parité~?
\item Que dit une parité aussi faible sur le simulateur~?
\item Quelles latences de production le teneur de cotation a-t-il subies, et que suppose le rejeu de recherche~?
\end{enumerate}

\medskip\noindent\textbf{Partie II --- P\&L.}
\begin{enumerate}[resume]
\item Donner le P\&L par heure à chaque étape de la décomposition.
\item Quelle est l'ampleur de chaque étape~?
\item Quelle part de l'étape de marché le hasard peut-il expliquer~?
\item Quelles étapes sont connues à l'avance, et lesquelles ont dû être mesurées~?
\end{enumerate}

\medskip\noindent\textbf{Partie III --- L'écart.}
\begin{enumerate}[resume]
\item Quelle part des lots de production provient d'ordres agressifs exécutés seulement contre le teneur de cotation, et pourquoi le rejeu ne peut-il pas les voir~?
\item Quelle part des lots de production a été exécutée sans autre ordre à leur prix, et quels étaient leurs mark-outs par rapport à ceux des autres~?
\item Pourquoi le teneur de cotation se retrouve-t-il seul à un prix~?
\item Qu'est-ce qui réduirait ces exécutions~?
\end{enumerate}

\medskip\noindent\textbf{Partie IV --- Le verdict.}
\begin{enumerate}[resume]
\item Énoncer le \emph{résultat nommé}~: la décomposition de l'écart entre le P\&L de production et le P\&L simulé en latence, marché, exécutions et frais.
\item Qu'apporte la calibration de la latence~?
\item Qu'est-ce que le passage au modèle de tête de file fait aux exécutions et au P\&L~?
\item Quelle est la correction structurelle~?
\item Que doit contenir le rapport quotidien de rapprochement~?
\item Comment utiliser l'\omterm{def:rs:simulation-versus-live:calibration}{implementation shortfall} pour les débouclages d'inventaire de cette stratégie~?
\item Que change-t-il si la stratégie trade 1~\% du volume à ses prix au lieu de 10~\%~?
\item En une phrase~: qu'est-ce qui fait d'un simulateur un instrument plutôt qu'une hypothèse~?
\end{enumerate}
\end{problem}

\section{Questions d'entretien}

\begin{interviewq}[$\star$ \iqroles{chercheur, trader}]\label{iq:rs:simulation-versus-live:1}
Votre stratégie fait en production la moitié de ce qu'elle faisait en simulation. Où regardez-vous en premier~?
\end{interviewq}

\begin{interviewq}[$\star\star$ \iqroles{chercheur}]\label{iq:rs:simulation-versus-live:2}
Qu'est-ce que l'\omterm{def:rs:simulation-versus-live:calibration}{implementation shortfall}, et comment la décomposez-vous~?
\end{interviewq}

\begin{interviewq}[$\star\star$ \iqroles{développeur, chercheur}]\label{iq:rs:simulation-versus-live:3}
Décrivez un \omterm{def:rs:simulation-versus-live:parity}{test de parité du rejeu} et ce qu'il détecte.
\end{interviewq}

\begin{interviewq}[$\star\star$ \iqroles{chercheur, trader}]\label{iq:rs:simulation-versus-live:4}
Pourquoi un \omterm{def:rs:vectorised-backtests:backtest}{backtest} sur des données de marché enregistrées pendant que vous tradiez fausse-t-il vos exécutions~?
\end{interviewq}

\begin{interviewq}[$\star\star$ \iqroles{chercheur}]\label{iq:rs:simulation-versus-live:5}
Comment calibreriez-vous un simulateur d'exécutions sur les exécutions de production sans le surajuster~?
\end{interviewq}

\begin{interviewq}[$\star\star\star$ \iqroles{chercheur, développeur}]\label{iq:rs:simulation-versus-live:6}
Concevez un rapprochement quotidien simulation--production pour une firme disposant de cinquante stratégies.
\end{interviewq}
